% ═══════════════════════════════════════════════════════════════════════════════
% PREGUNTAS Y OBJETIVOS - INSTITUTO ATENEO DE COLIMA
% Estudio Documental: Corrupción en México (2006-2012)
% ═══════════════════════════════════════════════════════════════════════════════

\nota{
  Este capítulo presenta la pregunta central y los objetivos de investigación 
  formulados conforme a los estándares de Hernández Sampieri (2014) para 
  estudios documentales de casos múltiples. La pregunta central guía el análisis 
  documental de los tres casos principales.
}

\section{Preguntas de investigación}

\subsection{Pregunta central}

¿De qué manera los esquemas de corrupción gestados durante el sexenio de Felipe Calderón (2006--2012) impactaron la integridad institucional, las finanzas públicas y los derechos humanos en México, a partir de los casos de Genaro García Luna, la Refinería Bicentenario y la Guardería ABC?

\subsection{Preguntas específicas}

\begin{enumerate}
  \item ¿De qué manera las distintas tipologías de corrupción ---desde prácticas incidentales hasta esquemas sistémicos y de captura estatal--- incidieron de forma verificable sobre el erario público, la legitimidad institucional y el ejercicio de los derechos humanos en México?
  
  \item ¿Cuál fue el nivel y la evolución empírica de la corrupción en México desde el cierre del sexenio 2006--2012 hasta las mediciones nacionales e internacionales más recientes (IPC 2025 y ENCIG 2025)?
  
  \item ¿Qué condiciones institucionales y mecanismos estructurales posibilitaron la captura del aparato de seguridad en el caso Genaro García Luna, los pasivos financieros en la Refinería Bicentenario y el desorden regulatorio con conflicto de interés en la Guardería ABC?
  
  \item ¿En qué medida difiere la magnitud verificable del daño ocasionado a la población y a sus derechos humanos respecto del tratamiento que reciben determinados escándalos bajo el encuadre mediático (framing), al contrastar los casos del periodo 2006--2012 con episodios contemporáneos como Veracruz y Dos Bocas?
  
  \item ¿Qué controles institucionales, mecanismos de rendición de cuentas y rediseños de incentivos sugiere la evidencia documental y teórica para reducir la brecha entre la normatividad anticorrupción y su aplicación efectiva en México?
\end{enumerate}

\section{Objetivos}

\subsection{Objetivo General}

Analizar los mecanismos causales y los efectos estructurales ---institucionales, financieros y en materia de derechos humanos--- de la corrupción en México durante el sexenio de Felipe Calderón (2006--2012), mediante un estudio documental de casos múltiples centrado en Genaro García Luna, la Refinería Bicentenario y la Guardería ABC, contrastando su impacto verificable con el tratamiento mediático del fenómeno.

\subsection{Objetivos Específicos}

\begin{enumerate}
  \item Sistematizar los fundamentos conceptuales y teóricos sobre la corrupción, sus tipologías jurídicas y sus efectos verificables sobre el erario, la legitimidad democrática y los derechos humanos, integrando doctrina internacional y mexicana.
  
  \item Describir el nivel, el costo económico y la trayectoria empírica de la corrupción en México a partir de indicadores secundarios oficiales de percepción agregada (IPC) y victimización directa (ENCIG), situando el cierre de 2012 en perspectiva longitudinal hasta 2026.
  
  \item Examinar el origen, los mecanismos de operación, las fallas de control y las resoluciones judiciales o de fiscalización de los tres casos emblemáticos del periodo 2006--2012 (Genaro García Luna, Refinería Bicentenario y Guardería ABC), junto con expedientes comparables del mismo sexenio.
  
  \item Contrastar el alcance real de las afectaciones a derechos humanos frente a la distorsión producida por el encuadre mediático (framing) mediante dos casos de contraste contemporáneo (Veracruz y Dos Bocas), derivando metainferencias sobre el fortalecimiento de la rendición de cuentas.
\end{enumerate}